\documentclass[../../main.tex]{subfiles}
\graphicspath{{../../image/}{../../figure/}}
\begin{document}

\section{带线性激发的平衡律}\label{section:ch11-带线性激发的平衡律}

我们考虑平衡律
\begin{align}
\partial_t u(x,t)+\partial_x f(u(x,t))-u(x,t)=0,\label{align:ch11-11-10-1}
\end{align}
其中 $f''(u)>0$, $-\infty<u<\infty$。为方便起见, 归一化 $f$ 及时空标架使 $f(0)=0$ 且 $f'(0)=0$。

目的是证明: 由于源项的去稳定作用与真正非线性的阻尼效应相互竞争, 零均值周期解渐近地成为有限振幅的驻波。

下文 $u(x,t)$ 表示\eqref{align:ch11-11-10-1}在上半平面上的可容许解, 局部有界变差, 初值 $u(x,0)=u_0(x)$ 在 $(-\infty,\infty)$ 上局部有界变差。

对无激波特征线, 系统\eqref{align:ch11-11-9-8}此处取形式
\begin{align}
\begin{cases}
\dot{\xi}=f'(v)\\
\dot{v}=v.
\end{cases}\label{align:ch11-11-10-2}
\end{align}
特别地, 分界线是在 $[0,\infty)$ 上无激波的特征线。显然, 沿分界线 $u$ 指数增长, 唯一例外是静止分界线 $x=\bar x$, 沿其上 $u$ 消失。如下命题识别静止分界线的发源点, 应与\cref{chapter:ch11-一维标量守恒律}比较。


\begin{lemma}\label{lemma:ch11-11-10-1}
直线 $x=\bar x$ 是与解 $u$ 关联的静止分界线, 当且仅当
\begin{align}
\int_{\bar x}^{z}u_0(x)\,\mathrm{d}x\geqslant0,\qquad -\infty<z<\infty,\label{align:ch11-11-10-3}
\end{align}
即 $\bar x$ 是 $u_0(\cdot)$ 的原函数的极小点。
\end{lemma}

\begin{proof}
之所以能如此精确地定位分界线发源点, 是因为齐次平衡律\eqref{align:ch11-11-10-1}可同样好地视为非齐次守恒律:
\begin{align}
\partial_t\left[e^{-t}u(x,t)\right]+\partial_x\left[e^{-t}f(u(x,t))\right]=0.\label{align:ch11-11-10-4}
\end{align}

先假设\eqref{align:ch11-11-10-3}成立。固定任意 $\bar t>0$, 考虑从点 $(\bar x,\bar t)$ 出发的极小后向特征线 $\xi(\cdot)$。在由 $\xi(\cdot)$ 图像、直线 $x=\bar x$ 与 $x$ 轴所围集合上积分\eqref{align:ch11-11-10-4}。应用 Green 公式, 并用\cref{theorem:ch11-11-9-3}与\eqref{align:ch11-11-10-2}, 得
\begin{align}
\int_0^{\bar t}e^{-t}\left\{f'(u(\xi(t),t))u(\xi(t),t)-f(u(\xi(t),t))\right\}\,\mathrm{d}t\nonumber\\
+\int_0^{\bar t}e^{-t}f(u(\bar x\pm,t))\,\mathrm{d}t+\int_{\bar x}^{z}u_0(x)\,\mathrm{d}x=0.\label{align:ch11-11-10-5}
\end{align}
\eqref{align:ch11-11-10-5}左端三项均非负, 故均须消失。于是 $u(\bar x\pm,t)=0$, $t\in(0,\infty)$, 从而 $x=\bar x$ 是静止分界线。

反之, 假设 $x=\bar x$ 是静止分界线, 沿其上 $u$ 消失。固定任意 $z<\bar x$。对 $\varepsilon>0$, 设 $\chi(\cdot)$ 是从点 $(z,0)$ 出发、斜率 $\dot{\chi}(t)=f'(\varepsilon e^t)$ 的曲线。设 $\chi(\cdot)$ 在时刻 $\bar t$ 与直线 $x=\bar x$ 相交。在集合 $\{(x,t):0\leqslant t\leqslant\bar t,\ \chi(t)\leqslant x\leqslant\bar x\}$ 上积分\eqref{align:ch11-11-10-4}并应用 Green 公式。加上并减去仅依赖 $t$ 的项, 最终得
\begin{align}
\int_0^{\bar t}e^{-t}\left\{f(\varepsilon e^t)-f(u(\chi(t)+,t))-f'(\varepsilon e^t)\left[\varepsilon e^t-u(\chi(t)+,t)\right]\right\}\,\mathrm{d}t\nonumber\\
-\int_0^{\bar t}e^{-t}f(\varepsilon e^t)\,\mathrm{d}t=\int_z^{\bar x}[u_0(x)-\varepsilon]\,\mathrm{d}x.\label{align:ch11-11-10-6}
\end{align}
上式左端两项均非正, 故右端也非正。令 $\varepsilon\downarrow0$, 得\eqref{align:ch11-11-10-3}, 对任意 $z<\bar x$。$z>\bar x$ 情形由同样方法处理。证毕。
\end{proof}

接下来证明, 相邻静止分界线之间, 解渐近达到有限振幅的驻波剖面。如下命题应与\cref{theorem:ch11-11-7-2}比较。

\begin{lemma}\label{lemma:ch11-11-10-2}
假设 $x=x_-$ 与 $x=x_+$, $x_-<x_+$, 是与解 $u$ 关联的相邻分界线, 即\eqref{align:ch11-11-10-3}对 $\bar x=x_-$ 与 $\bar x=x_+$ 成立, 但对区间 $(x_-,x_+)$ 内任何 $\bar x$ 不成立。考虑从点 $\left(\frac{x_-+x_+}{2},0\right)$ 出发的任意前向特征线 $\psi(\cdot)$。则当 $t\to\infty$ 时,
\begin{align}
u(x\pm,t)=
\begin{cases}
v_-(x)+o(1), & \text{在 }x_-<x<\psi(t)\text{ 上一致},\\
v_+(x)+o(1), & \text{在 }\psi(t)<x<x_+\text{ 上一致},
\end{cases}\label{align:ch11-11-10-7}
\end{align}
其中 $v_-(x)$ 与 $v_+(x)$ 是微分方程 $\partial_x f(v)=v$ 的解, 满足 $v_-(x_-)=0$ 与 $v_+(x_+)=0$。此外,
\begin{align}
\psi(t)=x_0+o(1),\label{align:ch11-11-10-8}
\end{align}
其中 $x_0$ 由条件
\begin{align}
\int_{x_-}^{x_0}v_-(y)\,\mathrm{d}y+\int_{x_0}^{x_+}v_+(y)\,\mathrm{d}y=0.\label{align:ch11-11-10-9}
\end{align}
确定。

特别地, 若 $u_0$ 在 $x_\pm$ 处可微且 $u_0'(x_\pm)>0$, 则\eqref{align:ch11-11-10-7}与\eqref{align:ch11-11-10-8}中的 $o(1)$ 阶升级为指数阶: $O(e^{-t})$。
\end{lemma}

\begin{proof}
当 $t\to\infty$ 时, 从点 $(\psi(t),t)$ 出发的极小后向特征线 $\zeta(\cdot)$ 收敛到被囚禁在区间 $[x_-,x_+]$ 内的分界线, 因而必是静止的。由于 $x_-$ 与 $x_+$ 相邻, $\zeta(\cdot)$ 必收敛到 $x_-$。特别地, $\zeta(0)=x_-+o(1)$, $t\to\infty$。

固定 $t>0$, 任取 $x\in(x_-,\psi(t)]$。设 $\xi(\cdot)$ 是从点 $(x,t)$ 出发的极小后向特征线, 在 $x$ 轴上被截于 $\xi(0)=\xi_0$, 其中 $x_-\leqslant\xi_0\leqslant\zeta(0)$。特别地, $\xi_0=x_-+o(1)$, $t\to\infty$。回忆\cref{theorem:ch11-11-9-3}, 积分系统\eqref{align:ch11-11-10-2}得 $v(\tau)=\bar u e^\tau$, $0\leqslant\tau\leqslant t$, 其中 $\bar u$ 满足 $u_0(\xi_0-)\leqslant\bar u\leqslant u_0(\xi_0+)$, 且
\begin{align}
x-\xi_0=\int_0^t f'(v(\tau))\,\mathrm{d}\tau=\int_{\bar u}^{u(x-,t)}\frac{f'(v)}{v}\,\mathrm{d}v.\label{align:ch11-11-10-10}
\end{align}

由\eqref{align:ch11-11-10-10}, $u(x-,t)=O(1)$, 从而 $\bar u=e^{-t}u(x-,t)=O(e^{-t})$。进而由 $\xi_0=x_-+o(1)$ 与 $\bar u=O(e^{-t})$, \eqref{align:ch11-11-10-10}给出\eqref{align:ch11-11-10-7}的上半段。当 $u_0'(x_-)>0$ 时, $\bar u=O(e^{-t})$ 蕴含 $\xi_0=x_-+O(e^{-t})$, 故 $o(1)$ 升级为 $O(e^{-t})$。\eqref{align:ch11-11-10-7}的下半段由同样方法处理。

在 $[x_-,x_+]\times[0,t]$ 上积分\eqref{align:ch11-11-10-1}, 推出 $\displaystyle\int_{x_-}^{x_+}u(x,t)\,\mathrm{d}x=0$, 于是由\eqref{align:ch11-11-10-7}得\eqref{align:ch11-11-10-8}与\eqref{align:ch11-11-10-9}。证毕。
\end{proof}

当初值 $u_0(\cdot)$ 是均值为 $M$ 的周期函数时, 解 $u(\cdot,t)$ 在时刻 $t$ 也是周期的, 均值为 $Me^t$, 因而当 $t\to\infty$ 时爆破, 除非 $M=0$。若 $M=0$, 每个周期区间内至少有一个 $\bar x$ 使\eqref{align:ch11-11-10-3}成立。因此,\cref{lemma:ch11-11-10-2}有如下推论, 与\cref{theorem:ch11-11-7-3}类似。

\begin{theorem}\label{theorem:ch11-11-10-3}
当初值 $u_0$ 是零均值周期函数时, 当 $t\to\infty$, 解 $u$ 趋于由形如\eqref{align:ch11-11-10-7}的小波构成的周期锯齿剖面。每个周期内小波的个数等于任意周期区间内使\eqref{align:ch11-11-10-3}成立的点 $\bar x$ 的个数。在一般情形, 即\eqref{align:ch11-11-10-3}在每个周期区间内仅在单点 $\bar x$ 处成立时, $u$ 趋于每个周期单齿的锯齿剖面。
\end{theorem}

\end{document}
